% !TeX root = ../papers/nc-p3.tex
\section{両側二重色の疎な台の全被覆}
\label{sec:klein-two-doublet-supports}
\begin{proposition}\label{prop:klein-two-doublet-supports}
二切断の色階数が $(2,2,1,1)$ の置換である正則表現型四元群窓を取る。
少なくとも一つの階数1長方形に零の行または列があるとする。
標準Koszul左右AS窓の台の候補は、厚い色の方向が平行なとき148軌道、異なるとき51軌道の表に尽くされる。
表は \path{computations/klein_doublet_parallel_masks.txt} と
\path{computations/klein_doublet_transverse_masks.txt} である。
\end{proposition}
\begin{proof}
二つの厚い色集合を $\{0,1\},\{0,q\}$ に平行移動し、平行なら $q=1$、
異なれば共通の線型変換で $q=2$ とする。各切断で残る二つの長方形にはそれぞれ $225^2$ 通りの選択がある。
二切断で指定される語の重なりは16セルであり、この上で一致する組を全て結合する。
48セルの指定後に未定なのは16セルである。

同時文字平行移動を含む $(t,t+h,t+h+k,t+k)$ の加算 ($h=0,1$, $k=0,q$)、
巡回、反転を作用させる。$q=2$ の奇数回の巡回に文字1,2の交換を伴わせ、
$q=1$ では文字2,3の同時交換も許す。これらは階数指定を保存する。
まず長方形の組の軌道代表を一つだけ選び、各代表について16セルの全ての補完を調べる。
元の台の個数を数えるときは、長方形の軌道の大きさを掛ける。全台の軌道に直すときは、
代表を保つ群の作用で補完をまとめる。二切断で指定された48成分が全て非零の軌道は、補完の前に除く。
\[
\begin{array}{c|rr}
 & q=1&q=2\\\hline
重なりが一致する長方形の組 & 1033395 & 973921\\
その軌道数（全非零も含む） &4856&8005\\
疎な補完で四つの階数2の台条件を通過&922893330&828689160\\
自由辺の強連結成分に分岐のない台&476008&292056\\
その軌道数&2779&2794\\
連続純関係・一文字商を通過する軌道&148&51
\end{array}
\]
階数2の台条件には零成分に対応する行直線と列の直交直線の一致を用い、
自由辺・純関係・一文字商の検査は片側欠色の場合と同じである。
各軌道は $16(i+j)+4i+k$ のビット順で最小の非負整数として記録する。
\path{computations/verify_klein_two_doublet_supports.py} はC++17列挙で上段の数と正規代表を求め、
下段の全代表について三次の成長検査を行い、二表と項ごとに比較する。
\end{proof}
